\chapter{Sistemas de Timoshenko}
\label{ch:timoshenko}

\textcolor{red}{UNA PEQUEÑA INTRODUCCIÓN}

\section{Modelos básicos de vigas}
\label{ch:timoshenko:section:primerosmodelos}

Para dar inicio a esta primera sección, haremos un \textit{tour} por la teoría básica de vigas. Esto cumplirá dos propósitos: el primero será el de contextualizar el sistema de Timoshenko que presentaremos al final de la sección (dado que su origen viene motivado por trabajos previos sobre la dinámica de vigas), y el segundo será el de esclarecer (ligeramente) la física que se oculta tras las ecuaciones. Continuando con lo anterior, no se pretende dar una deducción física extensa, y se asume familiaridad de parte del lector con los conceptos básicos de mecánica de medios continuos. Existen numerosos textos capaces de proveer al lector interesado de dichos conocimientos, como por ejemplo \textcolor{red}{[CITA]}.

Con todo lo anterior, pasaremos ahora a deducir la ecuación de vigas de Euler-Bernoulli.

\subsection{Ecuación de vigas de Euler-Bernoulli}
\label{ch:timoshenko:section:primerosmodelos:subsection:eulerbernoulli}

Comenzamos entonces enunciando las hipótesis del modelo:

\begin{enumerate}
    \item Una \textit{viga} es un prisma recto cuya longitud es mayor que sus otras dimensiones; así, esta teoría es especialmente adecuada para vigas esbeltas.

    \textcolor{red}{DIBUJO}

    \item La viga está hecha de un material \textit{linealmente elástico}. Es decir, obedece la ley de Hooke

    \item Las secciones transversales de la viga son \textit{simétricas} con respecto a los planos verticales\footnotemark. Luego, la \textit{línea} o \textit{eje neutro} (que separa, imaginariamente, la zona comprimida de la traccionada) corta a estas secciones y es perpendicular a ellas.

    \footnotetext{La simetría hace que el producto de inercia sea cero, eliminando el acoplamiento entre flexiones en distintos planos.}
    \textcolor{red}{DIBUJO}

    \item Los planos perpendiculares a la línea neutra permanecen planos y perpendiculares tras la deformación. Es decir, no se produce una \textit{deformación por cortante}.

    \item Se desprecian los efectos de la \textit{inercia rotacional}.

    \textcolor{red}{DIBUJO}

    \item Se supone que el ángulo de rotación es pequeño, por lo que se adopta la hipótesis de pequeñas \textit{deformaciones}.

    \textcolor{red}{DIBUJO}

    \item El material de la viga es homogéneo e isótropo (sus propiedades físicas son iguales en todas las direcciones), con una densidad $\rho$.
\end{enumerate}

\textcolor{red}{DIBUJO}

Ahora fijemos algo de notación. Dado que en esta teoría despreciamos las deformaciones por cortante, los efectos de la inercia rotatoria de las secciones y además asumimos que las secciones transversales permanecen planas y perpendiculares al eje neutro deformado, entonces se tiene que la respuesta de la viga ante una \textit{carga} queda descrita únicamente por la flexión. 

De esta manera, si denotamos por $(X,Y,Z)$ a un punto sobre la viga (a lo que nos referiremos como sus \textit{coordenadas materiales}), y suponemos que la longitud de la viga se extiende por el eje $X$, su altura por el $Y$ y su anchura por el $Z$, lo anterior nos permite colapsar la información de cada sección transversal al eje neutro (el cual se sitúa en la coordenada $Y=0$), por lo que todas nuestras funciones dependerán exclusivamente de la coordenada material $X$. 

Nuestro interés será entonces modelizar la deformación del eje neutro, y en consecuencia definimos como $\varphi(X, t)$ al \textit{desplazamiento} de dicho eje en la posición $X$, con $X \in [0,L]$ siendo $L$ la longitud de la viga, en el instante de tiempo $t$. Consideraremos que la función $\varphi(X,t)$ es negativa cuando la deformación se produce en el sentido positivo del eje $Y$, y positiva en el caso contrario. 

Por último, modelizaremos la carga mencionada anteriormente como una función $q(X)$, que representa la fuerza distribuida por unidad de longitud, y la cual puede hacer referencia tanto al peso de la propia viga, como a fuerzas externas que actúen sobre la misma.

\textcolor{red}{DIBUJO}

Hagamos ahora un estudio del estado de equilibrio de la viga. Es decir, una vez ocurrida la deformación, cuando las fuerzas se han equilibrado. Veamos entonces un diagrama de la viga deformada:

\textcolor{red}{DIBUJO}

Dado que haremos un estudio de la deformación una vez alcanzado el equilibrio, vamos a considerar la función $\varphi: [0,L] \rightarrow \mathbb{R}$, que será la que modelice la curva que define el eje neutro deformado.

En la figura anterior, tenemos que $\varrho$ es el radio de curvatura de la curva $x \mapsto \Big(x, \varphi(x)\Big)$. Así:

\[\frac{1}{\varrho} = \frac{\varphi^{\prime \prime}(x)}{\Big(\sqrt{1 + \big(\varphi^{\prime}(x) \big)^2} \Big)^3} \approx \varphi^{\prime \prime}(x),\]

donde hemos hecho uso de la hipótesis de desplazamientos pequeños para obtener una aproximación de $\frac{1}{\varrho}$.

Recordemos ahora, con algo más de formalidad, el concepto de \textit{deformación}. En su versión más general, se entiende por deformación al cambio de un cuerpo respecto a una configuración inicial. Así, en este caso, la deformación será el cambio respecto a una longitud de arco inicial (recuérdese que, como consecuencia de nuestras hipótesis, el eje neutro no cambia su longitud. Sin embargo, las fibras en un entorno del mismo si). Entonces, si la longitud de arco inicial de una fibra es

\[L_0 = (\varrho - Y)d\theta \ \text{ (la misma que la del eje neutro)},\]

y la longitud de arco de la fibra deformado es

\[L = \varrho d\theta,\]

entonces la deformación $\varepsilon$ será

\[\varepsilon = \frac{L - L_0}{L_0} = \frac{\varrho d\theta - (\varrho - Y)d\theta}{\varrho d\theta} = \frac{Y}{\varrho}.\]

Este resultado es coherente con la discusión llevada a cabo hasta el momento: En primer lugar, cuando $Y = 0$, obtenemos $\varepsilon = 0$; es decir, el eje neutro no sufre un cambio de longitud, su longitud permanece la misma. Por otro lado, ese cambio relativo a una posición inicial se hace (linealmente) más ``brusco'' en el resto de fibras según mayor sea la distancia vertical (en valor absoluto) al eje neutro.

Revisemos ahora el concepto de \textit{tensión}. Se define tensión (a veces también llamada \textit{estrés} o \textit{esfuerzo}) a la representación de la fuerza por unidad de área en el entorno de un punto material sobre una superficie de un medio continuo. Se produce a consecuencia de las fuerzas exteriores.

De manera intuitiva, conviene entonces pensar en las tensiones de un material como fuerzas internas que surgen dentro del mismo por la acción de fuerzas externas. Esta intuición es entonces congruente con la tercera ley de Newton.

Luego, como (por hipótesis) estamos ante un material elástico lineal, con deformaciones pequeñas, homogéneo e isótropo, entonces la ley de Hooke nos indica que la tensión viene dada por:

\[\sigma = E \varepsilon = E \frac{Y}{\varrho},\]

donde $E$ es una constante (para materiales linealmente elásticos) denominada \textit{módulo de Young}, la cual caracteriza la rigidez de un material frente a la flexión. Como una pequeña nota técnica, se recuerda que tanto la tensión como la deformación no son, en general, magnitudes escalares, sino magnitudes tensoriales. En este caso, por nuestras hipótesis y el tipo de cargas que asumimos, la única tensión presente es la tensión normal en la dirección del eje $X$; es decir, $\sigma_{XX} = \sigma(Y) = E \frac{Y}{\varrho}$. En consecuencia, podemos hablar de tensión y deformación como una magnitud puramente escalar.

Calculemos ahora el momento $M$ con relación al eje de simetría producido por la fuerza de tensión:

\[M = \int_A \sigma Y \, dA = \int_A E \frac{Y^2}{\varrho} \, dA = \frac{E}{\varrho} \int_A Y^2 \, dA = \frac{E I}{\varrho},\]

donde $A$ es el área de una sección transversal\footnotemark e $I = \int_A Y^2 \, dA$ es el momento de inercia. Considerando ahora la aproximación hecha de $\frac{1}{\varrho}$, se tiene que:

\footnotetext{A lo largo de este texto, salvo que se indique lo contrario, estaremos asumiendo que el área de las secciones transversales se mantiene constante. Sin embargo, la teoría no es limitante en este aspecto. El único objetivo de considerar el área constante es el de simplificar las ecuaciones.}
\[\frac{E I}{\varrho} = EI \varphi^{\prime \prime}(X)\]

\textcolor{red}{DIBUJO}
\textcolor{red}{DIBUJO}

Si denotamos ahora por $S$ al esfuerzo cortante o de cizallamiento, entonces, llevando a cabo un equilibrio de fuerzas obtenemos:

$(M + dM) - M - (S + dS)dX - q dX (\frac{dX}{2}) = 0.$

A pesar de que no estamos considerando deformaciones por cizallamiento, por simple estática es necesario tomar en consideración estos esfuerzos. Luego, despreciando los productos de diferenciales:

\[dM - SdX = 0 \implies \frac{dM}{dX} = S.\]

Ahora planteamos el equilibrio de fuerzas verticales:

\[S - (S + dS) + q dX = 0 \implies \frac{dS}{dX} = q.\]

Juntando ambas relaciones y sustituyendo por la expresión hallada de $M$ se obtiene la \textit{ecuación en estática de Euler-Bernoulli}:

\begin{align*}
    \label{eq::eulerbernoulliestatica}
    EI \frac{d^4\varphi}{dX^4} &= q
\end{align*}

Para el caso de una viga oscilante, en la que no se alcanza el equilibrio de las fuerzas verticales, debemos aplicar la segunda ley de Newton. En este caso consideramos $\varphi(X,t)$ y un trozo de viga de longitud $\Delta X$, donde el área de su sección transversal es $A$. Así,

\begin{align*}
    \sum F_{\text{verticales}} &= m a\\
    \implies S - (S + dS) + q\Delta X &= \int_{X}^{X + \Delta X} A \rho \frac{\partial^2 \varphi}{\partial t^2} \, dX\\
    \implies -dS + q\Delta X &= \int_{X}^{X + \Delta X} A \rho \frac{\partial^2 \varphi}{\partial t^2} \, dX.
\end{align*}

Dividiendo por $\Delta X$ los dos extremos de la igualdad, y tomando el límite cuando $\Delta X \rightarrow 0$ se obtiene:

\[A \rho \frac{\partial^2 \varphi}{\partial t^2} = -\frac{dS}{dX} + q,\]

e incorporando la relación $\frac{dS}{dX} = EI \frac{\partial^4 \varphi}{\partial X^4}$ derivada anteriormente, obtenemos la \textit{ecuación de evolución de Euler-Bernoulli}:

\begin{align*}
    %\label{eq::eulerbernoullievolucion}
    A \rho \frac{\partial^2 \varphi}{\partial t^2} + EI \frac{\partial^4 \varphi}{\partial X^4} = q.
\end{align*}

\subsection{El sistema de Timoshenko}
\label{ch:timoshenko:section:primerosmodelos:subsection:timoshenko}

Como se ha ido asomando a lo largo de la sección anterior, la ecuación de vigas de Euler-Bernoulli es adecuada para describir el comportamiento de ciertas vigas esbeltas, hechas de un material lo suficientemente homogéneo y ``bien portado''. Sin embargo, las hipótesis asumidas son, quizás, demasiado simples para el comportamiento de algunos sistemas de vigas que se presentan en la realidad. Por ejemplo, cuando se trata de vigas cortas o profundas en las que la deformación por cortante se vuelve relevante, cuando se tienen vigas de materiales no elásticos o no lineales (por ejemplificar alguno, los materiales compuestos son anisotrópicos, lo que quiere decir que, como conjunto, no presentan un comportamiento lineal), cuando las secciones transversales se distorsionan con la deformación, cuando se estudian vibraciones de alta frecuencia, entre otros. 

Debido a la necesidad de describir vigas con comportamientos más complejos y realistas, el ingeniero ucraniano \textit{Sthephen P. Timoshenko} (1878-1972) postula un sistema de ecuaciones en derivadas parciales que describe la dinámica de una viga teniendo en consideración el desplazamiento vertical y la rotación de la sección transversal respecto al eje neutro, permitiendo que esta no permanezca necesariamente perpendicular al eje de la viga. El desarrollo original de este sistema se puede encontrar en los trabajos de Timoshenko \textcolor{red}{[CITA]}.

La viga planteada por Timoshenko conserva entonces algunas de las hipótesis planteadas para la viga de Euler-Bernoulli. Más explícitamente, la viga de Timoshenko está hecha de un material linealmente elástico, las secciones transversales son simétricas respecto a los planos verticales (lo que permite la existencia de un eje neutro), se asumen deformaciones pequeñas y el material de la viga es homogéneo e isótropo. Asimismo, se relajan o eliminan algunas de las hipótesis. El cambio más notable respecto a la viga de Euler-Bernoulli es que la viga de Timoshenko no asume que las secciones perpendiculares al eje neutro permanecen planas y perpendiculares tras la deformación. Luego, si se consideran los efectos de la inercia rotacional y se relaja la hipótesis de la longitud frente al resto de dimensiones de la viga (ahora se permiten vigas más cortas o profundas).

Ahora fijamos, una vez más, algo de notación: En una viga de longitud $L$, denotaremos al \textit{desplazamiento vertical} y al \textit{ángulo de rotación de una sección transversal respecto a la normal al eje neutro} por $\varphi = \varphi(X,t)$ y $\psi = \psi(X,t)$ respectivamente, dependiendo ambas aplicaciones de una posición $X \in [0, L]$ y un tiempo $t \geq 0$. Véase entonces la siguiente figura:

\textcolor{red}{DIBUJO}

Es importante recalcar que, respecto al modelo de Euler-Bernoulli, se ha introducido una nueva variable, que es la \textit{rotación de la sección transversal} $\psi$. Como se ha relajado la hipótesis que dicta que las secciones planas y normales al eje neutro permanecen planas y normales tras la deformación, entonces se sigue que $\psi \neq \frac{\partial \varphi}{\partial X}$. 

A partir de este momento, dejaremos de referirnos a las coordenadas materiales como $(X,Y,Z)$ y empezaremos a referirnos a ellas como $(x,y,z)$ con el fin de aligerar la notación. 

Siguiendo entonces con el desarrollo de Timoshenko \textcolor{red}{[CITA]}, se tiene que las \textit{ecuaciones de equilibrio dinámico} para un elemento diferencial de viga, que relacionan los esfuerzos internos con las aceleraciones traslacionales y rotacionales, se expresan como:

\begin{align}
    \rho A \varphi_{tt} &= S_x, \label{timoshenkoBasico:equilDinamicoCortante}\\
    \rho I \psi_{tt} &= M_x - S, \label{timoshenkoBasico:equilDinamicoFlector}
\end{align}

donde (al igual que en el modelo de la sección anterior) $\rho$ es la \textit{densidad de masa}, $A$\footnotemark e $I$ son el área y el momento de inercia de una sección transversal de la viga, respectivamente; $S$ es el esfuerzo de cizallamiento (que ahora, debido a las nuevas hipótesis, no es despreciable) y $M$ el momento flector.  

\footnotetext{Nuevamente, asumimos el área constante.}

Por otra parte, las \textit{relaciones constitutivas elásticas} para el esfuerzo cortante (de cizallamiento) y el momento flector vienen dadas por:

\begin{align}
    S &= k^{\prime}GA(\varphi_x + \psi), \label{timoshenkoBasico:relConstitutivasCortante}\\
    M &= EI \psi_x, \label{timoshenkoBasico:relConstitutivasFlector}
\end{align}

y se tiene que $k^{\prime}$ es el \textit{factor de corrección por cortante}, mientras que $G$ y $E$ son el \textit{módulo de elasticidad transversal} (o bien \textit{módulo de cizalladura/cortadura}) --que caracteriza la rigidez de un material frente a esfuerzos de cizallamiento-- y el \textit{módulo de Young}, respectivamente. 

El factor de corrección $k^{\prime}$ se introduce porque de las hipótesis planteadas anteriormente se deduce que el esfuerzo cortante es uniforme en toda la sección transversal (los desarrollos originales de Timoshenko \textcolor{red}{[CITA]} contemplan la deducción e implicaciones de este hecho. Dado que el interés de este texto no radica en la física, sólo hemos de mencionar los detalles más relevantes); sin embargo, esto no es cierto en general, dado que el esfuerzo por cizallamiento va decayendo desde el centro de la viga hacia los extremos. Por ejemplo, en vigas con secciones rectangulares se conoce que el esfuerzo por cortadura sigue una distribución parabólica. Entender de manera intuitiva el esfuerzo por cizallamiento como aquel que surge por la resistencia al deslizamiento de las capas (fibras) internas del material en respuesta a una fuerza puede resultar esclarecedor, pues la cara superior e inferior no tienen fibras por encima o por debajo de ellas que las ``empujen'' o ``frenen'', mientras que las fibras más centrales tienen todas las capas por encima y por debajo ``empujando'' y ``frenando'' (tracción y compresión).

\textcolor{red}{DIBUJO}

Sustituyendo \eqref{timoshenkoBasico:relConstitutivasCortante}-\eqref{timoshenkoBasico:relConstitutivasFlector} en \eqref{timoshenkoBasico:equilDinamicoCortante}-\eqref{timoshenkoBasico:equilDinamicoFlector} y abreviando las constantes como

\[\rho_1 = \rho A, \ \rho_2 = \rho I, \ k = k^{\prime}GA, \ b = EI,\]

obtenemos el \textit{sistema de vigas de Timoshenko}

\begin{equation}
\label{eqTimoshenkoBasico}
\begin{cases}
\rho_1 \varphi_{tt} - k(\varphi_x + \psi)_x = 0 \, &\text{en } (0,L) \times (0, \infty),\\
\rho_2 \psi_{tt} - b\psi_{xx} + k(\varphi_x + \psi) = 0 \, &\text{en } (0,L) \times (0, \infty).
\end{cases}
\end{equation}

Asumiremos entonces que la viga está fija en los extremos $\{0,L\}$, lo que nos lleva a considerar las siguientes \textit{condiciones de fronteras}:

\[\varphi(0,t) = \varphi(L,t) = 0, \ \psi(0,t) = \psi(L,t) = 0, \ t \geq 0.\]

Asimismo, para $t = 0$ consideramos las \textit{condiciones iniciales}:

\[\varphi(x,0) = \varphi_0(x), \ \varphi_t(x,0) = \varphi_1(x), \ \psi(x,0) = \psi_0(x), \ \psi_t(x,0) = \psi_1(x), \ x \in (0, L).\]

Lo anterior resulta entonces en el \textit{problema de valor inicial y de frontera}:

\begin{equation}
\label{pviTimoshenkoBasico}
\begin{cases}
\rho_1 \varphi_{tt} - k(\varphi_x + \psi)_x = 0 \, &\text{en } (0,L) \times (0, \infty),\\
\rho_2 \psi_{tt} - b\psi_{xx} + k(\varphi_x + \psi) = 0 \, &\text{en } (0,L) \times (0, \infty), \\
\varphi(x,0) = \varphi_0(x), \varphi_t(x,0) = \varphi_1(x) \, &\text{en } (0, L) \times \{0\}, \\
\psi(x,0) = \psi_0(x), \psi_t(x,0) = \psi_1(x) \, &\text{en } (0, L) \times \{0\}, \\
\varphi(0,t) = \varphi(L,t) = 0 \, &\text{en } \{0, L\} \times [0, \infty),\\
\psi(0,t) = \psi(L,t) = 0 \, &\text{en } \{0, L\} \times [0, \infty).
\end{cases}
\end{equation}

\subsection{La energía en el sistema de Timoshenko}
\label{ch:timoshenko:section:primerosmodelos:subsection:timoshenkoEnergia}

Al problema de valor inicial y de frontera \eqref{pviTimoshenkoBasico} visto anteriormente se le puede asociar el siguiente funcional de energía definido por:

\begin{equation}
    \label{timoshenkoBasico:funcionalEnergia}
    \begin{aligned}
        E(t)\text{\footnotemark} \ = \  &\frac{1}{2} \Bigg[ \rho_1 \int_0^L \big(\varphi_t(x,t) \big)^2 \, dx + \rho_2 \int_0^L \big(\psi_t(x,t) \big)^2 \, dx + b\int_0^L \big(\psi_x(x,t) \big)^2 \, dx\\
        &+ k\int_0^L \big(\varphi_x(x,t) + \psi(x,t) \big)^2 \, dx \Bigg], \quad t \geq 0
    \end{aligned}
\end{equation}

\footnotetext{\color{red} Dado el marco funcional en el que se estudiará este problema, todas las integrales presentadas en este funcional y en el resto del documento se interpretan en el sentido de Lebesgue.}

\textcolor{red}{PONER DEDUCCIÓN FÍSICA DEL FUNCIONAL}

Para poder describir el estado del sistema utilizaremos un vector $U$ definido de la siguiente manera:

\begin{align*}
    U(t) &= \Big(\varphi(\cdot, t), \varphi_t(\cdot, t), \psi(\cdot, t), \psi_t(\cdot, t)\Big)^T,\\
    U_t(t) &= \Big(\varphi_t(\cdot, t), \varphi_{tt}(\cdot, t), \psi_t(\cdot, t), \psi_{tt}(\cdot, t)\Big)^T,\\
    U(0) &= \Big(\varphi_0(\cdot), \varphi_1(\cdot), \psi_0(\cdot), \psi_1(\cdot)\Big)^T.
\end{align*}

La elección de este vector como estado del sistema no es arbitraria, y a lo largo de este texto se nos presentarán varias oportunidades para justificarla. Por adelantarnos a esto: tenemos que, desde una perspectiva física, el estado del sistema viene determinado por su configuración y sus velocidades iniciales (``funciones posición'' y sus derivadas). En efecto, en mecánica, conocer únicamente la posición no es suficiente para predecir la evolución, siendo necesario especificar también las velocidades.

Desde un punto de vista matemático (y aunque se ahondará más al respecto en el siguiente capítulo \textcolor{red}{[PONER CITA PROX CAP]}), ahora mismo bastará con decir que esta elección nos permitirá hacer una descripción de la dinámica del sistema apta para la maquinaría matemática que se pretende utilizar: la teoría de semigrupos de operadores lineales. 

Así pues, dado que el sistema de Timoshenko \eqref{pviTimoshenkoBasico} es de segundo orden en tiempo, los datos iniciales naturales vienen dados por las variables y sus derivadas temporales. La introducción de estas últimas permitirá reformular el sistema como una ecuación de evolución de primer orden en un espacio de Hilbert adecuado. Veremos además que esta elección es coherente con el funcional de energía asociado \eqref{timoshenkoBasico:funcionalEnergia}.

Habiendo asumido el vector $U$ como vector de estado, entonces, y debido a una necesidad física, hemos de imponer que

\[E(t) < \infty.\]

Es decir, debemos imponer que la energía exista y sea finita. Esto nos permite motivar la elección de un \textit{espacio de fase}, que es el espacio al que pertenecerá el vector $U$. Así, se ve implicado que:

\begin{align}
    &\int_0^L \big(\varphi_t(x,t) \big)^2 \, dx < \infty,\label{varphi_tL2}\\
    &\int_0^L \big(\psi_t(x,t) \big)^2 \, dx < \infty,\label{psi_tL2}\\
    &\int_0^L \big(\psi_x(x,t) \big)^2 \, dx < \infty,\label{psiH01}\\
    &\int_0^L \big(\varphi_x(x,t) + \psi(x,t) \big)^2 \, dx < \infty\label{varphiH01}.
\end{align}

En consecuencia, para cada $t \geq 0$, de \eqref{varphi_tL2} y \eqref{psi_tL2} se deduce que

\[
\varphi_t(\cdot,t)\in L^2(0,L), \quad \psi_t(\cdot,t)\in L^2(0,L).
\]

Además, de \eqref{psiH01} se obtiene que

\[
\psi_x(\cdot,t)\in L^2(0,L).
\]

Dado que estamos manejando una única dimensión espacial, el hecho de que $\psi_x \in L^2(0,L)$, junto con la condición de contorno $\psi(0,t) = 0$, nos permite hacer el siguiente desarrollo:

\begin{align*}
    \left (\int_0^L |\psi_x(x,t)| \, dx \right )^2 &= \left (\int_0^L |1 \cdot \psi_x(x,t)| \, dx \right )^2\\
    \text{(Cauchy–Schwarz)} &\leq \left (\int_0^L |\psi_x(x,t)|^2 \, dx \right ) \left (\int_0^L dx \right )\\
    &= L \int_0^L |\psi_x(x,t)|^2 \, dx < \infty.
\end{align*}

Asimismo, por \textcolor{green}{LEMA 8.2 BREZIS} se puede definir una función $v$:

\[v(x) = \int_0^x \psi_x(x,t) \, dx, \ x \in [0,L],\]

y $v \in C(0,L)$. Además, 

\[\int_0^L v \phi ' \, dx = - \int_0^L \psi_x \phi \, dx, \ \forall \phi \in C_0^\infty(0,L).\]

Luego, necesariamente $v = \psi$ y así se tiene entonces que $\psi \in C(0,L)$ y que $\psi(x,t) = \int_0^x \psi_x(x,t) \, dx $. 

Procedemos ahora a acotar la norma de $\psi$ de la siguiente manera:

\[
|\psi(x,t)|^2 = \left| \int_0^x \psi_x(s,t) \, ds \right|^2 \leq \left( \int_0^x 1^2 \, ds \right) \left( \int_0^x |\psi_x(s,t)|^2 \, ds \right) \leq x \int_0^L |\psi_x(s,t)|^2 \, ds.
\]

Integrando esta expresión en $x$ sobre el intervalo $(0,L)$, obtenemos

\[
\int_0^L |\psi(x,t)|^2 \, dx \leq \frac{L^2}{2} \int_0^L |\psi_x(x,t)|^2 \, dx < \infty.
\]

Esta es, de hecho, una demostración directa de la \textit{desigualdad de Poincaré} en dimensión uno. Por tanto, se concluye que $\psi(\cdot,t) \in L^2(0,L)$. Al tener $\psi$ y $\psi_x$ en $L^2(0,L)$ junto con las condiciones de contorno nulas, se sigue que
\[
\psi(\cdot,t)\in H_0^1(0,L).
\]
Por otro lado, de \eqref{varphiH01} se tiene que
\[
\varphi_x(\cdot,t)+\psi(\cdot,t)\in L^2(0,L).
\]
Como el espacio $L^2(0,L)$ es un espacio vectorial y ya hemos demostrado que $\psi(\cdot,t)\in L^2(0,L)$, deducimos inmediatamente que $\varphi_x(\cdot,t)\in L^2(0,L)$. De nuevo, valiéndonos de la condición de contorno $\varphi(0,t) = 0$ y repitiendo el argumento integral anterior, garantizamos que $\varphi(\cdot,t) \in L^2(0,L)$. Por consiguiente,
\[
\varphi(\cdot,t)\in H_0^1(0,L).
\]

Con esta caracterización, resulta natural definir el espacio de fase para el sistema de Timoshenko como el espacio de Hilbert\footnotemark:

\footnotetext{Una demostración de que la suma directa de espacios de Hilbert es también un espacio de Hilbert puede ser hayada en \textcolor{green}{CONWAY: A COURSE IN FUNCTIONAL ANALYSIS. CAP 1.6}}

\begin{align}
    \mathcal{H} = H_0^1(0,L) \times L^2(0,L) \times H_0^1(0,L) \times L^2(0,L),\label{espacioDeFaseTimoshenko}
\end{align}
de modo que, para cada $t \geq 0$, el vector de estado satisfaga $U(t) \in \mathcal{H}$.

Es importante hacer una precisión sobre el sentido en el que deben entenderse las derivadas espaciales (como $\varphi_x$ y $\psi_x$) que aparecen tanto en el funcional \eqref{timoshenkoBasico:funcionalEnergia} como en la definición del espacio de fase. 

Originalmente, el sistema \eqref{pviTimoshenkoBasico} se deduce bajo la suposición de que las variables de estado son lo suficientemente regulares como para admitir derivadas en el sentido clásico. Sin embargo, exigir este nivel de regularidad (soluciones clásicas) resulta ser demasiado restrictivo desde el punto de vista matemático, ya que el espacio de las funciones continuamente diferenciables no es completo bajo las normas integrales de espacios funcionales como $L^2$. Esto significa que el límite de una sucesión de funciones suaves podría dejar de ser diferenciable en el sentido clásico, impidiendo usar herramientas como la teoría de semigrupos.

Para sortear esta dificultad y garantizar la existencia de soluciones, es necesario relajar el concepto de derivada. Por lo tanto, en el marco de los espacios de Sobolev $H_0^1(0,L)$ y $L^2(0,L)$, todas las derivadas espaciales se entienden en el sentido débil o distribucional. Es decir, decimos que $\varphi_x \in L^2(0,L)$ es la derivada débil de $\varphi \in L^2(0,L)$ si satisface la fórmula de integración por partes:

\[
\int_0^L \varphi(x) \phi'(x) \, dx = - \int_0^L \varphi_x(x) \phi(x) \, dx, \quad \text{para toda función de prueba } \phi \in C_c^\infty(0,L).
\]

Esta formulación débil es precisamente la que fundamenta la estructura del espacio de Hilbert $\mathcal{H}$ y dota de sentido riguroso al estado del sistema, abarcando configuraciones que, aunque no sean derivables clásicamente en todo punto, son físicamente admisibles.

Dicho todo lo anterior, veamos ahora una de las propiedades más relevantes del modelo de vigas de Timoshenko. Propiedad que influenciará casi por completo lo que resta de este texto.\\

\begin{Proposicion}{(Conservación de la energía en la viga de Timoshenko)}
    \label{ch:timoshenko:prop:ConservacionEnergiaTimoshenkoBasico}

    El funcional de energía definido en \eqref{timoshenkoBasico:funcionalEnergia} verifica:

    \[\frac{d}{dt}E(t) = 0, \ \ \forall t > 0.\]
\end{Proposicion}

\begin{proof}
    En primer lugar, consideramos $\varphi, \varphi_t, \psi$ y $\psi_t$ en los espacios funcionales descritos anteriormente. Ahora, empezamos multiplicando la primera ecuación del sistema \eqref{pviTimoshenkoBasico} por $\varphi_t$ y la segunda por $\psi_t$, e integramos en $(0,L)$:

    \begin{equation*}
    \begin{cases}
    \displaystyle \rho_1 \int_0^L \varphi_{tt} \varphi_t \, dx - k \int_0^L (\varphi_x + \psi)_x \varphi_t \, dx = 0\\
    \displaystyle \rho_2 \int_0^L \psi_{tt} \psi_t \, dx - b\int_0^L \psi_{xx}\psi_t \, dx + k \int_0^L (\varphi_x + \psi) \psi_t \, dx = 0
    \end{cases}
    \end{equation*}

    Ahora, realizamos integración por partes en los términos $\int_0^L (\varphi_x + \psi)_x \varphi_t \, dx$ y $\int_0^L \psi_{xx}\psi_t \, dx$:

    \begin{align*}
        &\int_0^L (\varphi_x + \psi)_x \varphi_t \, dx = \cancelto{0}{\varphi_t(\varphi_x + \psi)\Big|_0^L} - \int_0^L(\varphi_x + \psi) \varphi_{tx} \, dx,\\
        &\int_0^L \psi_{xx}\psi_t \, dx = \cancelto{0}{\psi_t \psi_x\Big|_0^L} - \int_0^L \psi_x \psi_{tx} \, dx.
    \end{align*}

    Sustituyendo lo anterior y sumando ambas ecuaciones:

    \[\rho_1 \int_0^L\varphi_{tt}\varphi_t \, dx + \rho_2 \int_0^L \psi_{tt}\psi_t \, dx + b\int_0^L \psi_x \psi_{tx} \, dx + k \int_0^L (\varphi_x + \psi)(\varphi_{tx} + \psi_t) \, dx = 0.\]

    Considerando ahora las igualdades

    \begin{align*}
        (\varphi_x + \psi)(\varphi_{tx} + \psi_t) &= \frac{1}{2}\frac{\partial}{\partial t}(\varphi_x + \psi)^2,\\
        \varphi_{tt}\varphi_t &= \frac{1}{2}\frac{\partial}{\partial t}(\varphi_t)^2,\\
        \psi_{tt}\psi_t &= \frac{1}{2}\frac{\partial}{\partial t}(\psi_t)^2,\\
        \psi_x \psi_{tx} &= \frac{1}{2}\frac{\partial}{\partial t}(\psi_x)^2,
    \end{align*}

    donde en la primera igualdad hemos utilizado la conmutatividad de las derivadas distribucionales (``Teorema de Schwarz'').

    Sustituimos estas relaciones en la expresión hallada anteriormente y obtenemos:

    \begin{align*}
        &\frac{d}{dt}\frac{1}{2}\left [\rho_1 \int_0^L(\varphi_t)^2 \, dx + \rho_2 \int_0^L (\psi_t)^2 \, dx + b\int_0^L (\psi_x)^2 \, dx + k \int_0^L (\varphi_x + \psi)^2 \, dx \right] = 0\\
        \iff&\frac{d}{dt}E(t) = 0,
    \end{align*}

    que es exactamente lo que se quería demostrar.
\end{proof}